\scanpage{224}
% archive sheet number: 223

\mspage{- 7 -} \hfill n\textsuperscript{o}370

si $f\leq\sum\limits_{i=1}^n c_i\gamma(s_i)g$ avec des $c_i\geq0$, on a $\sup f\leq\sum\limits_{i=1}^n c_i\,g(s_i^{-1}s)$ pour un $s\in G$, donc $\sup f\leq(\sum\limits_{i=1}^n c_i)\sup g$, d'où (iv). Prouvons (v) ; soient $f\in\mathcal{K}$, $g,h\in\mathcal{K}_+^*$ ; si $f\leq\sum\limits_{i=1}^n c_i\gamma(s_i)g$ et $g\leq\sum\limits_{j=1}^p d_j\gamma(t_j)h$ ($c_i,d_j\geq0$, $s_i,t_j\in G$), $f\leq\sum\limits_{i,j}c_i\,d_j\,\gamma(s_it_j)h$, donc $(f:h)\leq\sum\limits_{i,j}c_i\,d_j=(\sum\limits_i c_i)(\sum\limits_j d_j)$ ; donc $(f:h)\leq(f:g)(g:h)$. Si on applique (v) à $f_{\unsure{o|0}}$, $f$, $g$ d'une part et à $f$, $f_o$, $g$ d'autre part, on obtient (vi). Enfin, soient $f,f',h\in\mathcal{K}_+$ avec $h(s)\geq1$ sur le support de $f+f'$, et soit $\varepsilon>0$. Posons $F=f+f'+\frac{1}{2}\varepsilon h$ ; les fonctions $\varphi,\varphi'$, qui coïncident respectivement avec $f/F$ et $f'/F$ sur le support de $f+f'$ et qui sont nulles en dehors de celui-ci, appartiennent à $\mathcal{K}_+$ ; pour tout $\eta>0$, il existe un voisinage compact $V$ de $e$ tel que $|\varphi(s)-\varphi(t)|\leq\eta$ et $|\varphi'(s)-\varphi'(t)|\leq\eta$ pour $s^{-1}t\in V$. Soit alors $g\in\mathcal{K}_+^*(V)$ ; pour tout $s\in G$, on a $\varphi.\,\gamma(s)g\leq(\varphi(s)+\eta).\,\gamma(s)g$ (en effet, c'est évident aux points où $\gamma(s)g$ s'annule, donc hors de $sV$ ; et, dans $sV$, $\varphi\leq\varphi(s)+\eta$) ; de même, $\varphi'.\,\gamma(s)g\leq(\varphi'(s)+\eta).\,\gamma(s)g$. Ceci posé, soient $c_1,\ldots,c_n$ des nombres $\geq0$ et $s_1,\ldots,s_n$ des éléments de $G$ tels que $F\leq\sum\limits_{i=1}^n c_i\gamma(s_i)g$ ; on a

\[
f=\varphi F\leq\sum_{i=1}^n c_i\,\varphi.\,\gamma(s_i)g\leq\sum_{i=1}^n c_i(\varphi(s_i)+\eta).\,\gamma(s_i)g
\]

et de même pour $f'$ ; par suite

\[
(f:g)+(f':g)\leq\sum_{i=1}^n c_i(\varphi(s_i)+\varphi'(s_i)+2\eta)\leq(1+2\eta)\sum_{i=1}^n c_i
\]

puisque $\varphi+\varphi'\leq1$. En appliquant la définition de $F$, puis (ii), (iii), et (v), on en conclut

\[
\begin{aligned}
(f:g)+(f':g)&\leq(1+2\eta)(F:g)\leq(1+2\eta)\left[(f+f':g)+\frac{1}{2}\varepsilon(h:g)\right]\leq{}\\
&\leq(f+f':g)+\frac{1}{2}\varepsilon(h:g)+2\eta(f+f':h)(h:g)+\varepsilon\eta(h:g)
\end{aligned}
\]

et, si on a choisi $\eta$ tel que $\eta\left[2(f+f':h)+\varepsilon\right]\leq\frac{1}{2}\varepsilon$, on obtient (vii).

Quand $V$ parcourt l'ensemble des voisinages compacts de $e$, les $\mathcal{K}_+^*(V)$ forment une base de filtre $\unsure{\mathscr{B}|\mathcal{B}}$ sur $\mathcal{K}_+^*$. Soit $\unsure{\mathscr{F}|\mathcal{F}}$ un ultrafiltre sur $\mathcal{K}_+^*$ plus fin que $\mathscr{B}$.
